% Options for packages loaded elsewhere
\PassOptionsToPackage{unicode}{hyperref}
\PassOptionsToPackage{hyphens}{url}
\PassOptionsToPackage{dvipsnames,svgnames,x11names}{xcolor}
\documentclass[
  11pt,
]{article}
\usepackage{xcolor}
\usepackage[margin=2.3cm]{geometry}
\usepackage{amsmath,amssymb}
\setcounter{secnumdepth}{-\maxdimen} % remove section numbering
\usepackage{iftex}
\ifPDFTeX
  \usepackage[T1]{fontenc}
  \usepackage[utf8]{inputenc}
  \usepackage{textcomp} % provide euro and other symbols
\else % if luatex or xetex
  \usepackage{unicode-math} % this also loads fontspec
  \defaultfontfeatures{Scale=MatchLowercase}
  \defaultfontfeatures[\rmfamily]{Ligatures=TeX,Scale=1}
\fi
\usepackage{lmodern}
\ifPDFTeX\else
  % xetex/luatex font selection
\fi
% Use upquote if available, for straight quotes in verbatim environments
\IfFileExists{upquote.sty}{\usepackage{upquote}}{}
\IfFileExists{microtype.sty}{% use microtype if available
  \usepackage[]{microtype}
  \UseMicrotypeSet[protrusion]{basicmath} % disable protrusion for tt fonts
}{}
\makeatletter
\@ifundefined{KOMAClassName}{% if non-KOMA class
  \IfFileExists{parskip.sty}{%
    \usepackage{parskip}
  }{% else
    \setlength{\parindent}{0pt}
    \setlength{\parskip}{6pt plus 2pt minus 1pt}}
}{% if KOMA class
  \KOMAoptions{parskip=half}}
\makeatother
% definitions for citeproc citations
\NewDocumentCommand\citeproctext{}{}
\NewDocumentCommand\citeproc{mm}{%
  \begingroup\def\citeproctext{#2}\cite{#1}\endgroup}
\makeatletter
 % allow citations to break across lines
 \let\@cite@ofmt\@firstofone
 % avoid brackets around text for \cite:
 \def\@biblabel#1{}
 \def\@cite#1#2{{#1\if@tempswa , #2\fi}}
\makeatother
\newlength{\cslhangindent}
\setlength{\cslhangindent}{1.5em}
\newlength{\csllabelwidth}
\setlength{\csllabelwidth}{3em}
\newenvironment{CSLReferences}[2] % #1 hanging-indent, #2 entry-spacing
 {\begin{list}{}{%
  \setlength{\itemindent}{0pt}
  \setlength{\leftmargin}{0pt}
  \setlength{\parsep}{0pt}
  % turn on hanging indent if param 1 is 1
  \ifodd #1
   \setlength{\leftmargin}{\cslhangindent}
   \setlength{\itemindent}{-1\cslhangindent}
  \fi
  % set entry spacing
  \setlength{\itemsep}{#2\baselineskip}}}
 {\end{list}}
\usepackage{calc}
\newcommand{\CSLBlock}[1]{\hfill\break\parbox[t]{\linewidth}{\strut\ignorespaces#1\strut}}
\newcommand{\CSLLeftMargin}[1]{\parbox[t]{\csllabelwidth}{\strut#1\strut}}
\newcommand{\CSLRightInline}[1]{\parbox[t]{\linewidth - \csllabelwidth}{\strut#1\strut}}
\newcommand{\CSLIndent}[1]{\hspace{\cslhangindent}#1}
\setlength{\emergencystretch}{3em} % prevent overfull lines
\providecommand{\tightlist}{%
  \setlength{\itemsep}{0pt}\setlength{\parskip}{0pt}}
\usepackage{bookmark}
\IfFileExists{xurl.sty}{\usepackage{xurl}}{} % add URL line breaks if available
\urlstyle{same}
% fallback for those not using the hyperref driver hyperxmp:
\makeatletter
\@ifundefined{xmpquote}{\newcommand{\xmpquote}[1]{#1}}{}
\makeatother
\hypersetup{
  colorlinks=true,
  linkcolor={Maroon},
  filecolor={Maroon},
  citecolor={Blue},
  urlcolor={Blue},
  pdfcreator={LaTeX via pandoc}}

\author{}
\date{}

\begin{document}

\section{INTRODUCTION\_BACKGROUND.md --- bullet bank for §1
(merge-ready)}\label{introduction_background.md-bullet-bank-for-1-merge-ready}

Purpose: background material for the Introduction, kept as
\textbf{bullets} so you can merge/trim them yourself. Every bullet
carries its \texttt{@bibkey} (all keys are in
\texttt{manuscript/references.bib} and were verified against publisher
pages on 23 Nov 2026 --- see \texttt{LITERATURE\_SURVEY.md}). Bullets
marked \textbf{{[}needs citation{]}} assert something true but not yet
backed by a verified reference.

\textbf{Web of Science status:} the WoS interface loads
\textbf{IP-authenticated for your institute} (the footer shows PANDIT
DEENDAYAL PETROLEUM UNIV), but its Nextgen Angular search form ignores
synthetic input from automation --- only real keystrokes register. Two
ways to bring WoS results in: (i) type one query in the visible tab
(e.g.~\texttt{TS=("low\ flow"\ AND\ CMIP6\ AND\ India)}) and tell me,
and I will harvest the result list; or (ii) paste the WoS search URL
after you run it. Meanwhile the bullets below use the publisher-verified
set plus \textbf{OpenAlex} searches run 23 Nov 2026 (verified metadata
and citation counts), which surfaced the India-propagation and
ANOVA-canon references now added (Bhardwaj 2020, Gupta 2024, Bosshard
2013, Chen 2011, Clark 2016).

\begin{center}\rule{0.5\linewidth}{0.5pt}\end{center}

\subsection{A. Why low flows decide water security in Peninsular
India}\label{a.-why-low-flows-decide-water-security-in-peninsular-india}

\begin{itemize}
\tightlist
\item
  Peninsular India's water availability is decided in the non-monsoon
  months: irrigation, drinking-water supply, hydropower and
  environmental flows all draw on low flows, while \textasciitilde80--90
  \% of annual rainfall arrives in the monsoon. \textbf{{[}needs
  citation: monsoon fraction{]}}
\item
  Meteorological and agricultural drought are already pervasive and well
  mapped: the Drought Atlas of India provides 0.05° monitoring for
  1901--2020 (Chuphal et al. 2024) --- but it is not a
  streamflow-drought product.
\item
  Peninsular catchments are dominated by hard-rock aquifers with limited
  deep storage, so rainfall deficits translate quickly into streamflow
  deficits; propagation from meteorological to hydrological drought is
  controlled by catchment storage, baseflow index and fractured-rock
  fraction (Barker et al. 2016).
\item
  Consequence for modelling: the storage terms of a conceptual model ---
  not the monsoon flush --- set the low-flow response, which is why a
  four-parameter structure (GR4J) and its six-parameter low-flow variant
  (GR6J) can agree on mean flow while disagreeing on droughts
  (Pushpalatha et al. 2011; Perrin et al. 2003).
\item
  Abstraction and regulation are not modelled here; they are flagged as
  a limitation and deferred to the anthropogenic-drought phase.
  \textbf{{[}needs citation: CAMELS-IND regulation attributes{]}}
\end{itemize}

\subsection{B. The projection landscape over India (what
exists)}\label{b.-the-projection-landscape-over-india-what-exists}

\begin{itemize}
\tightlist
\item
  Bias-corrected CMIP6 at 0.25° for South Asia (Mishra et al. 2020):
  broad-scale, not catchment-tuned, and not hydrologically evaluated at
  low flows.
\item
  India-wide streamflow projections with four hydrological structures
  (MILC/GR4J/HBV/HYMOD), 12 GCMs, 128 basins (Sharma et al. 2026) ---
  the closest prior art; they note that ``uncertainty existed for low
  flow events'' without formalising it.
\item
  Observation-constrained projections report robust streamflow
  \textbf{increases} in Indian rivers (Chuphal and Mishra 2026) --- a
  large-scale statement that says nothing about low-flow behaviour in
  dry catchments.
\item
  CAMELS-IND: 472 catchments, daily forcings and streamflow, 211
  attributes, 1981--2014 ({Mangukiya et al.} 2025); CAMELS-IND-Proj:
  LSTM streamflow for 13 GCMs × 4 SSPs, 2015--2100, now published in
  \emph{Scientific Data} ({Singh et al.} 2026).
\item
  Large-sample hydrological benchmarking over Peninsular India: 47
  conceptual models on 159 watersheds, GR4J in the five-model sufficient
  set ({Mangukiya et al.} 2026) --- establishes that structure choice
  matters, but benchmarks \emph{observed} flow, not projections or low
  flows specifically.
\item
  Single-basin Indian studies select GCMs on forcing skill or ML
  ensembles before hydrological modelling (Krishna basin; Papagni
  (Naresh and Mathew 2026); Betwa) --- useful precedents, but
  single-basin and not low-flow-specific.
\end{itemize}

\subsection{C. The selection problem in the methods
literature}\label{c.-the-selection-problem-in-the-methods-literature}

\begin{itemize}
\tightlist
\item
  Forcing-based selection (rank GCMs on how well corrected
  rainfall/temperature match observations) is the default and is
  inherited from global assessments; it says nothing about discharge
  skill.
\item
  Whole-flow criteria (KGE/NSE on the full hydrograph) are dominated by
  the monsoon season, where climate models are already similar; low-flow
  behaviour is a small fraction of the variance they see.
\item
  GR4J-based low-flow projections exist (Q95, 37 Irish catchments)
  (Meresa et al. 2022), showing the structural sensitivity of Q95 to
  climate change --- but without a selection framework or a second
  structure.
\item
  LSTM surrogates are trained across many basins with static attributes
  and evaluated regionally (Kratzert et al. 2018, 2019); best practice
  explicitly forbids single-basin training (Kratzert et al. 2024), so a
  benchmark must be judged on its cross-model medians --- which is what
  our audit does.
\item
  Uncertainty decomposition is established: ANOVA over GCM × scenario ×
  hydrologic-model factors in Korea for \textbf{low flows and drought
  events} (GR4J/GR6J/SWAT; model selection explains 21.9--51.9 \% of
  low-flow variance) (Park et al. 2025), in Korea for calibration
  condition/length (3.9--9.8 \% of runoff, 2.7 \% of drought) (Kim and
  Chung 2026), and in the Mediterranean for GCM-RCM × bias-correction
  method (Senatore et al. 2022).
\end{itemize}

\subsection{D. Why bias-correction tail design deserves to be a
factor}\label{d.-why-bias-correction-tail-design-deserves-to-be-a-factor}

\begin{itemize}
\tightlist
\item
  Quantile-mapping choices materially alter projected extremes,
  quantiles and trends (Cannon et al. 2015); error-correction method
  choice changes the climate-change signal (Themeßl et al. 2012;
  Gudmundsson et al. 2012).
\item
  Practice splits into \emph{clamp at the training range} (constant)
  versus \emph{extend linearly beyond it} (linear) --- a design decision
  usually made silently and never tested against the response variable.
\item
  Our own audit is a concrete instance of the failure class Cannon
  documents: the pre-fix linear implementation extended both tails with
  the slope of the lowest CDF interval (a drizzle artefact), turning a
  60 mm/d future storm into 4.0×10⁹ mm/d and corrupting 9,985 of 13,068
  corrected series before the guard caught it.
\item
  After correction (per-tail robust slope + 1,000 mm/d ceiling), the
  tail policy turns out to matter ≤0.5 \% of low-flow projection
  variance (H4) --- a \emph{bounded} negative finding with a practical
  message: tail-design anxiety is warranted for flood-oriented work, not
  for low-flow selection.
\end{itemize}

\subsection{E. The structural axis (why GR4J vs GR6J, not GR4J vs ``any
model'')}\label{e.-the-structural-axis-why-gr4j-vs-gr6j-not-gr4j-vs-any-model}

\begin{itemize}
\tightlist
\item
  GR6J adds an exponential-store exchange branch explicitly designed to
  improve low-flow simulation (Pushpalatha et al. 2011) --- the axis is
  the literature's own low-flow variant, not an arbitrary pair.
\item
  Published evidence does \textbf{not} establish GR6J superiority: GR4J
  sometimes reproduces observed low flows more closely in prolonged dry
  periods; therefore the design quantifies \textbf{divergence}, never
  superiority (pre-registered in \texttt{GR6J\_00\_setup.Rmd}).
\item
  Identical-treatment rule: same inputs, windows, warm-up, optimizer,
  objective, screening and out-of-sample evaluation for both structures
  --- the only difference is the structure itself.
\end{itemize}

\subsection{F. Drought propagation as a selection
dimension}\label{f.-drought-propagation-as-a-selection-dimension}

\begin{itemize}
\tightlist
\item
  SPI→SRI lag via cross-correlation of standardised indices is standard
  (Barker et al. 2016), and propagation under climate change is an
  active literature (Ireland (Meresa et al. 2023); Yangtze ({Ma et al.}
  2023)).
\item
  \textbf{India-specific propagation already exists and must be cited:}
  Bhardwaj, Shah, Aadhar \& Mishra (2020, \emph{JGR-Atmospheres}, 229
  citations) map the propagation of meteorological to hydrological
  drought across India (Bhardwaj et al. 2020); Gupta, Jain \& Pandey
  (2024, \emph{Hydrol. Process.}) track changing propagation times in a
  semi-arid Indian basin (Gupta et al. 2024). \textbf{Differentiation:}
  both characterise propagation; neither turns reproduction of the
  observed propagation lag into a \textbf{GCM selection criterion} or
  asks whether mean-flow selection finds it --- that is our H3.
\item
  The \textbf{selection use} is what is new here: no prior study asks
  whether a GCM that reproduces the observed propagation lag is the one
  you should project droughts with.
\item
  Guardrail: propagation must be measured as the cross-correlation lag,
  never inferred from near-identity of SPI and SRI (the Machhar-type
  artefact); our diagnostic (mean \textbar SRI−SPI\textbar{} =
  0.29--0.34, far from the ≈0.013 signature) is reported as evidence.
\end{itemize}

\subsection{F2. The uncertainty-decomposition canon (for §4.2
positioning)}\label{f2.-the-uncertainty-decomposition-canon-for-4.2-positioning}

\begin{itemize}
\tightlist
\item
  ANOVA/variance decomposition of hydrological climate-impact ensembles
  has a canonical lineage: Chen, Brissette, Poulin \& Leconte (2011,
  \emph{WRR}, 469 citations) partitioned overall uncertainty for a
  Canadian watershed (Chen et al. 2011); \textbf{Bosshard et al.~(2013,
  \emph{WRR}, 389 citations) is the reference that formalised ANOVA over
  GCM × downscaling × hydrological-model factors} (Bosshard et al. 2013)
  --- cite it as the method's statistical anchor, alongside Park 2025
  and Kim \& Chung 2026 as the recent applications.
\item
  Clark et al.~(2016, \emph{Curr. Clim. Change Rep.}, 334 citations)
  frames the uncertainty source inventory (GCM structure, downscaling,
  bias correction, hydrological model structure, internal variability)
  ({Clark et al.} 2016) --- use it to motivate \emph{why} a three-factor
  decomposition with a tail-design axis is a defensible subset rather
  than an arbitrary one.
\item
  These three make the concession in §2 airtight: the \emph{method} is
  canonical, the \emph{factor set} (tail design) and the
  \emph{application} (India, low flows, 192 catchments) are the
  contribution.
\end{itemize}

\subsection{G. Ready-to-merge gap statements (one sentence
each)}\label{g.-ready-to-merge-gap-statements-one-sentence-each}

\begin{itemize}
\tightlist
\item
  No prior study tests whether forcing-skill rank and low-flow-skill
  rank diverge over Indian catchments --- the formal ρ test and its
  near-zero outcome are new here.
\item
  No prior Peninsular India study uses the GR4J/GR6J pair as a
  structural-uncertainty axis for low flows.
\item
  No prior decomposition treats bias-correction \textbf{tail design} as
  a first-order factor alongside GCM and hydrological structure, and
  none runs at 192-catchment scale over India with a low-flow metric
  layer.
\item
  Drought-propagation reproduction has been measured, but never used as
  a GCM selection criterion over Indian catchments.
\item
  CAMELS-IND-Proj has not yet been audited at low flows by an
  independent process-based user; the dry-tail question is open in the
  dataset paper's own framing.
\end{itemize}

\subsection{H. Numbers ready to drop into the text (all
executed)}\label{h.-numbers-ready-to-drop-into-the-text-all-executed}

\begin{itemize}
\tightlist
\item
  H1: ρ = −0.027, 95\% CI {[}−0.095, 0.038{]}, 1,000 bootstrap
  resamples, seed 20261118; 78 non-ephemeral catchments; pre-registered
  failure threshold ρ \textgreater{} 0.6.
\item
  H2: dry-tercile 7-day minimum divergence −449 pp; Q95 −111 pp;
  whole-flow logNSE immaterial (median −0.01 to 0.05, p = 0.41--0.45);
  172 paired catchments.
\item
  H3: propagation share 0.91--0.96 for all 22 models; spread \textless{}
  0.10; ρ(forcing rank, propagation share) = −0.350; 167 frozen observed
  lags.
\item
  H4: structure ≥ 5\% in 149 of 180 summary cells; dry-tercile structure
  share 0.66 / 0.68 / 0.59 (Q95 / 7-day minimum / low-tail shape) vs GCM
  0.23 / 0.23 / 0.21; mean flow GCM 0.95--0.96; tail policy ≤ 0.005;
  11,520 cells (11,514 decomposed, 6 logged).
\item
  H5: LSTM win shares 0.86 / 0.88 logNSE, 0.92 / 0.95 low-flow volume,
  0.81 / 0.78 7-day minimum, 0.71 / 0.71 Q95, 0.39 / 0.39 dry spell (vs
  GR4J / GR6J).
\item
  Roster: 472 (setting) → 242 (Tier-1/2) → 192 (both structures;
  64/64/64) → 172 (H2 base) / 78 (H1 composite) / 167 (propagation
  lags); GCM coverage 14/13/13/14.
\end{itemize}

\subsection{I. Citation items still open (for you, at
camera-ready)}\label{i.-citation-items-still-open-for-you-at-camera-ready}

\begin{itemize}
\tightlist
\item
  Monsoon rainfall fraction and low-flow water-use statistics for
  Peninsular India.
\item
  Anthropogenic abstraction/regulation quantification (CAMELS-IND
  attributes; the Year-4 phase).
\item
  Full citations: Meresa 2023 (done), Ma 2023 author list, Pushpalatha
  2011 end page, Krishna-basin GCM-selection paper authors, Mishra 2020
  article number.
\item
  WoS/Scopus coverage and citation counts under your institutional
  login.
\end{itemize}

\protect\phantomsection\label{refs}
\begin{CSLReferences}{1}{1}
\bibitem[\citeproctext]{ref-barker2016propagation}
Barker, Lucy J., Jamie Hannaford, Alison Chiverton, and Cecilia
Svensson. 2016. {``From Meteorological to Hydrological Drought Using
Standardised Indicators.''} \emph{Hydrology and Earth System Sciences}
20: 2483--505. \url{https://doi.org/10.5194/hess-20-2483-2016}.

\bibitem[\citeproctext]{ref-bhardwaj2020propagation}
Bhardwaj, Kunal, Deep Shah, Saran Aadhar, and Vimal Mishra. 2020.
{``Propagation of Meteorological to Hydrological Droughts in {India}.''}
\emph{Journal of Geophysical Research: Atmospheres} 125: e2020JD033455.
\url{https://doi.org/10.1029/2020JD033455}.

\bibitem[\citeproctext]{ref-bosshard2013uncertainty}
Bosshard, Thomas, Maria Carambia, Klaus Goergen, Sven Kotlarski, Stefan
Krähenmann, and Massimiliano Zappa. 2013. {``Quantifying Uncertainty
Sources in an Ensemble of Hydrological Climate-Impact Projections.''}
\emph{Water Resources Research} 49: 1523--36.
\url{https://doi.org/10.1029/2011WR011533}.

\bibitem[\citeproctext]{ref-cannon2015qm}
Cannon, Alex J., Stephen R. Sobie, and Trevor Q. Murdock. 2015. {``Bias
Correction of {GCM} Precipitation by Quantile Mapping: How Well Do
Methods Preserve Changes in Quantiles and Extremes?''} \emph{Journal of
Climate} 28 (17): 6938--59.
\url{https://doi.org/10.1175/JCLI-D-14-00754.1}.

\bibitem[\citeproctext]{ref-chen2011uncertainty}
Chen, Jie, François P. Brissette, Annie Poulin, and Robert Leconte.
2011. {``Overall Uncertainty Study of the Hydrological Impacts of
Climate Change for a {Canadian} Watershed.''} \emph{Water Resources
Research} 47: W12542. \url{https://doi.org/10.1029/2011WR010602}.

\bibitem[\citeproctext]{ref-chuphal2024droughtatlas}
Chuphal, Dharm Singh, Apurv Prakash Kushwaha, Saran Aadhar, and Vimal
Mishra. 2024. {``Drought Atlas of {India}, 1901--2020.''}
\emph{Scientific Data} 11: 7.
\url{https://doi.org/10.1038/s41597-023-02862-0}.

\bibitem[\citeproctext]{ref-chuphal2026obsconstrained}
Chuphal, Dharm Singh, and Vimal Mishra. 2026. {``Observation-Constrained
Projections Reveal Robust Streamflow Increases in {Indian} Rivers.''}
\emph{Earth's Future} 14.

\bibitem[\citeproctext]{ref-clark2016uncertainty}
{Clark, Martyn P., Robert L. Wilby, Ethan D. Gutmann, et al.} 2016.
{``Characterizing Uncertainty of the Hydrologic Impacts of Climate
Change.''} \emph{Current Climate Change Reports} 2: 119--28.
\url{https://doi.org/10.1007/s40641-016-0034-x}.

\bibitem[\citeproctext]{ref-gudmundsson2012eqm}
Gudmundsson, Luc, J. B. Bremnes, J. E. Haugen, and T. Engen-Skaugen.
2012. {``Technical Note: Downscaling {RCM} Precipitation to the Station
Scale Using Statistical Transformations -- a Comparison of Methods.''}
\emph{Hydrology and Earth System Sciences} 16: 3383--90.
\url{https://doi.org/10.5194/hess-16-3383-2012}.

\bibitem[\citeproctext]{ref-gupta2024propagation}
Gupta, Ajay, Manoj Kumar Jain, and Rajendra Prasad Pandey. 2024. {``The
Changing Characteristics of Propagation Time from Meteorological Drought
to Hydrological Drought in a Semi-Arid River Basin in {India}.''}
\emph{Hydrological Processes}, ahead of print.
\url{https://doi.org/10.1002/hyp.15266}.

\bibitem[\citeproctext]{ref-kimchung2026hess}
Kim, J. H., and E.-S. Chung. 2026. {``Insights into Uncertainties in
Future Drought Analysis Using Hydrological Simulation Model.''}
\emph{Hydrology and Earth System Sciences} 30: 227--47.
\url{https://doi.org/10.5194/hess-30-227-2026}.

\bibitem[\citeproctext]{ref-kratzert2024nevertrain}
Kratzert, Frederik, Martin Gauch, Daniel Klotz, and Grey Nearing. 2024.
{``{HESS Opinions}: Never Train a {Long Short-Term Memory} ({LSTM})
Network on a Single Basin.''} \emph{Hydrology and Earth System Sciences}
28: 4187--201. \url{https://doi.org/10.5194/hess-28-4187-2024}.

\bibitem[\citeproctext]{ref-kratzert2018lstm}
Kratzert, Frederik, Daniel Klotz, Christoph Brenner, Karsten Schulz, and
Mathias Herrnegger. 2018. {``Rainfall--Runoff Modelling Using {Long
Short-Term Memory} ({LSTM}) Networks.''} \emph{Hydrology and Earth
System Sciences} 22: 6005--20.
\url{https://doi.org/10.5194/hess-22-6005-2018}.

\bibitem[\citeproctext]{ref-kratzert2019regional}
Kratzert, Frederik, Daniel Klotz, Sepp Hochreiter, and Grey S. Nearing.
2019. {``Towards Learning Universal, Regional, and Local Hydrological
Behaviors via Machine Learning Applied to Large-Sample Datasets.''}
\emph{Hydrology and Earth System Sciences} 23: 5089--111.
\url{https://doi.org/10.5194/hess-23-5089-2019}.

\bibitem[\citeproctext]{ref-ma2023yangtze}
{Ma, F. et al.} 2023. {``Intensification of Drought Propagation over the
{Yangtze River} Basin Under Climate Warming.''} \emph{International
Journal of Climatology} 43 (12): 5640--58.
\url{https://doi.org/10.1002/joc.8165}.

\bibitem[\citeproctext]{ref-mangukiya2026wrr}
{Mangukiya, Nishant Kumar et al.} 2026. {``Benchmarking and Selecting
Optimal Hydrological Models for Catchment-Scale Streamflow Simulation in
{Peninsular India}.''} \emph{Water Resources Research} 62.

\bibitem[\citeproctext]{ref-mangukiya2025camelsind}
{Mangukiya, Nishant Kumar, Anurag Sharma, Ankita Wany, et al.} 2025.
{``{CAMELS-IND}: Hydrometeorological Time Series and Catchment
Attributes for 472 Catchments in Peninsular India.''} \emph{Earth System
Science Data} 17: 461--91.
\url{https://doi.org/10.5194/essd-17-461-2025}.

\bibitem[\citeproctext]{ref-meresa2022water}
Meresa, Hadush, Shaun Donegan, Samaneh Golian, and Conor Murphy. 2022.
{``Simulated Changes in Seasonal and Low Flows with Climate Change for
{Irish} Catchments.''} \emph{Water} 14 (10): 1556.
\url{https://doi.org/10.3390/w14101556}.

\bibitem[\citeproctext]{ref-meresa2023jgr}
Meresa, Hadush, Conor Murphy, and Shaun Donegan. 2023. {``Propagation
and Characteristics of Hydrometeorological Drought Under Changing
Climate in {Irish} Catchments.''} \emph{Journal of Geophysical Research:
Atmospheres} 128: e2022JD038025.
\url{https://doi.org/10.1029/2022JD038025}.

\bibitem[\citeproctext]{ref-mishra2020biascorrected}
Mishra, Vimal, Udit Bhatia, and Aadhar Tiwari. 2020. {``Bias-Corrected
Climate Projections for {South Asia} from {CMIP6}.''} \emph{Scientific
Data} 7: 1--13. \url{https://doi.org/10.1038/s41597-020-00681-1}.

\bibitem[\citeproctext]{ref-naresh2026papagni}
Naresh, C., and A. Mathew. 2026. {``Machine-Learning-Enhanced
Multi-Model Climate Ensemble for Future Rainfall, Temperature and
Drought Assessment in the Semi-Arid {Papagni River} Basin, {India}.''}
\emph{Journal of Atmospheric and Solar-Terrestrial Physics}.

\bibitem[\citeproctext]{ref-park2025jhrs}
Park, J., J. H. Kim, J. H. Sung, and E. S. Chung. 2025. {``Quantifying
Uncertainty Contributions of Hydrologic Modeling Process to Hydrologic
Drought Projection in {South Korea}.''} \emph{Journal of Hydrology:
Regional Studies} 62: 102991.
\url{https://doi.org/10.1016/j.ejrh.2025.102991}.

\bibitem[\citeproctext]{ref-perrin2003gr4j}
Perrin, Charles, Claude Michel, and Vazken Andréassian. 2003.
{``Improvement of a Parsimonious Model for Streamflow Simulation.''}
\emph{Journal of Hydrology} 279: 275--89.
\url{https://doi.org/10.1016/S0022-1694(03)00225-7}.

\bibitem[\citeproctext]{ref-pushpalatha2011gr6j}
Pushpalatha, Raji, Charles Perrin, Nicolas Le Moine, and Vazken
Andréassian. 2011. {``A Downward Structural Sensitivity Analysis of
Hydrological Models to Improve Low-Flow Simulation.''} \emph{Journal of
Hydrology} 411: 66--79.
\url{https://doi.org/10.1016/j.jhydrol.2011.10.022}.

\bibitem[\citeproctext]{ref-senatore2022med}
Senatore, Alfonso, Domenico Fuoco, Mario Maiolo, Giuseppe Mendicino,
Gerhard Smiatek, and Harald Kunstmann. 2022. {``Evaluating the
Uncertainty of Climate Model Structure and Bias Correction on the
Hydrological Impact of Projected Climate Change in a {Mediterranean}
Catchment.''} \emph{Journal of Hydrology: Regional Studies} 42: 101120.
\url{https://doi.org/10.1016/j.ejrh.2022.101120}.

\bibitem[\citeproctext]{ref-sharma2026hp}
Sharma, V., S. Barbhuiya, and V. Gupta. 2026. {``Projections of Future
Streamflow for {India} Informed by {CMIP6} Global Climate Models.''}
\emph{Hydrological Processes} 40 (3): e70476.

\bibitem[\citeproctext]{ref-singh2026scidata}
{Singh, A., N. K. Mangukiya, G. Sahu, et al.} 2026.
{``{CAMELS-IND-Proj}: Daily Streamflow Projections for 472 {Peninsular
Indian} Catchments Under Future Climate Scenarios.''} \emph{Scientific
Data}, ahead of print. \url{https://doi.org/10.1038/s41597-026-08858-w}.

\bibitem[\citeproctext]{ref-themessl2012ecc}
Themeßl, Matthias Jakob, Andreas Gobiet, and Georg Heinrich. 2012.
{``Empirical-Statistical Downscaling and Error Correction of Regional
Climate Models and Its Impact on the Climate Change Signal.''}
\emph{Climatic Change} 112 (2): 449--68.
\url{https://doi.org/10.1007/s10584-011-0224-4}.

\end{CSLReferences}

\end{document}
